\documentclass[
10pt, % Main document font size
a4paper, % Paper type, use 'letterpaper' for US Letter paper
oneside, % One page layout (no page indentation)
%twoside, % Two page layout (page indentation for binding and different headers)
headinclude,footinclude, % Extra spacing for the header and footer
BCOR5mm, % Binding correction
]{scrartcl}
%%%%%%%%%%%%%%%%%%%%%%%%%%%%%%%%%%%%%%%%%
% Arsclassica Article
% Structure Specification File
%
% This file has been downloaded from:
% http://www.LaTeXTemplates.com
%
% Original author:
% Lorenzo Pantieri (http://www.lorenzopantieri.net) with extensive modifications by:
% Vel (vel@latextemplates.com)
%
% License:
% CC BY-NC-SA 3.0 (http://creativecommons.org/licenses/by-nc-sa/3.0/)
%
%%%%%%%%%%%%%%%%%%%%%%%%%%%%%%%%%%%%%%%%%

%----------------------------------------------------------------------------------------
%	REQUIRED PACKAGES
%----------------------------------------------------------------------------------------

\usepackage[
nochapters, % Turn off chapters since this is an article        
beramono, % Use the Bera Mono font for monospaced text (\texttt)
eulermath,% Use the Euler font for mathematics
pdfspacing, % Makes use of pdftex’ letter spacing capabilities via the microtype package
dottedtoc % Dotted lines leading to the page numbers in the table of contents
]{classicthesis} % The layout is based on the Classic Thesis style

\usepackage{arsclassica} % Modifies the Classic Thesis package
\usepackage{ulem}
\usepackage{graphicx}
\usepackage{float}
\usepackage{algorithm2e}
\usepackage[export]{adjustbox}
\usepackage[utf8]{inputenc}  
\usepackage[english,frenchb]{babel} 
\usepackage[T1]{fontenc} % Required for accented characters
\usepackage{multicol}
\usepackage{graphicx} % Required for including images
\graphicspath{{Figures/}} % Set the default folder for images

\usepackage{enumitem} % Required for manipulating the whitespace between and within lists

\usepackage{lipsum} % Used for inserting dummy 'Lorem ipsum' text into the template


\expandafter\def\csname ver@subfig.sty\endcsname{}


\usepackage{svg}
\usepackage{hyperref}

\usepackage{amsmath,amssymb,amsthm} % For including math equations, theorems, symbols, etc

\usepackage{varioref} % More descriptive referencing

%----------------------------------------------------------------------------------------
%	THEOREM STYLES
%---------------------------------------------------------------------------------------

\theoremstyle{definition} % Define theorem styles here based on the definition style (used for definitions and examples)
\newtheorem{definition}{Definition}

\newtheorem*{remark}{Remark}
\theoremstyle{plain} % Define theorem styles here based on the plain style (used for theorems, lemmas, propositions)
\newtheorem{theorem}{Theorème}
\newtheorem{lemma}[theorem]{Lemme}
\newtheorem{proposition}[theorem]{Proposition}

\theoremstyle{remark} % Define theorem styles here based on the remark style (used for remarks and notes)

%----------------------------------------------------------------------------------------
%	HYPERLINKS
%---------------------------------------------------------------------------------------

\hypersetup{
%draft, % Uncomment to remove all links (useful for printing in black and white)
colorlinks=true, breaklinks=true, bookmarks=true,bookmarksnumbered,
urlcolor=webbrown, linkcolor=RoyalBlue, citecolor=webgreen, % Link colors
pdftitle={}, % PDF title
pdfauthor={\textcopyright}, % PDF Author
pdfsubject={}, % PDF Subject
pdfkeywords={}, % PDF Keywords
pdfcreator={pdfLaTeX}, % PDF Creator
pdfproducer={LaTeX with hyperref and ClassicThesis} % PDF producer
}

\hyphenation{Fortran hy-phen-ation} 

\title{\normalfont{FTML Exercices 3}\\
} 
\author{\spacedlowsmallcaps{}
} 
\date{\today}
\begin{document}
\renewcommand{\sectionmark}[1]{\markright{\spacedlowsmallcaps{#1}}}
\lehead{\mbox{\llap{\small\thepage\kern1em\color{halfgray} \vline}\color{halfgray}\hspace{0.5em}\rightmark\hfil}} 
\pagestyle{scrheadings}
\maketitle 
\setcounter{tocdepth}{3} 

\tableofcontents

\section{\large\color{Blue}Estimateurs et espérances}

\textbf{Remarque:} attention, le terme estimateur est utilisé avec deux sens différents ici:
    \begin{itemize}
        \item Estimateur statistique

            \url{https://fr.wikipedia.org/wiki/Estimateur_(statistique)} 
        
        \item Estimateur au sens "modèle": fonction utilisée pour prédire en fonction d'une entrée $x$.
    \end{itemize}

\subsection{Moyenne empirique et espérance}

1) On se donne une variable aléatoire réelle X ayant un moment d’ordre 1. On note E[X] son espérance. On suppose que l’on dispose de n samples indépendants tirés de cette variable, $(x_1,\dots,x_n)$, qui suivent donc la loi de $X$.
On considère leur moyenne empirique $S_n =  \frac{1}{n}\sum_{i=1}^{n} x_i$. Montrer que l’espérance de $S_n$ est l'espérance de $X$, $E[X]$ (autrement dit, que $S_n$ est un estimateur non biaisé de $E[X]$).
\\

2) Simuler une variable de votre choix pour observer la convergence de $S_n$ vers $E[X]$.

\subsection{Espérance du risque empirique d'un estimateur indépendant du dataset}

1) On se donne un problème d’apprentissage supervisé usuel et un estimateur $f$, fixé et indépendant du dataset. Montrer que l’espérance du risque empirique de $f$ est le risque réel de $f$. Autrement dit: le risque empirique de f est un estimateur non biaisé de son risque réel.

La notion d’espérance a bien un sens ici car il faut voir le dataset comme une variable aléatoire, ainsi le risque empirique de $f$ est bien une variable aléatoire.
\\


2) Simuler un problème de votre choix pour observer le résultat (c’est ce qui est fait dans simulations/lecture\textunderscore 3/empirical\textunderscore risk\textunderscore convergence)

\subsection{Espérance du risque empirique du minimiseur du risque empirique}

Montrer que le résultat précédent n'est plus vrai quand $f$ est le minimiseur du risque empirique, c'est-à-dire la fonction $f_n$ pour laquelle le risque empirique prend la valeur la plus faible. On note que maintenant, $f_n$ est une variable aléatoire, calculée à partir du dataset.

\textbf{Remarque:} vous pouvez montrer le résultat par le calcul mais également le vérifier en faisant des simulations.

\textbf{Remarque:} Ce fait est une autre façon de dire qu'il est possible d'overfitter.

\end{document}
